% =============================================================================
% Ready-to-paste LaTeX: the shared genre space and the zero-shot cross-corpus
% transfer experiment. Answers supervisor requests:
%   "test on Blockbuster without training on it, train on Eerola, and compare with
%    the results of the paper itself"
%   "the classifier has to have the same genre classes in both datasets to be fair"
%
% Numbers from experiments/cross_dataset/exp_zero_shot.py (results/zero_shot.json),
% seed 42. Ma et al. figures are quoted from their Table 4.
%
% Placement: a Results section of its own ("Cross-Corpus Generalisation"), after the
% in-domain results. The label \label{sec:zero_shot} is referenced from
% docs/latex/evaluation_protocol.tex.
%
% CITATION NOTE: uses ma2021computational (already in bib/library.bib) and
% eerola2011genrespecific (BibTeX entry in emotion_genre_relationship.tex).
% =============================================================================

\section{Cross-Corpus Generalisation}
\label{sec:zero_shot}

All results reported so far are cross-validated within a single corpus. That design can
establish that a representation carries genre information, but not that the information
is a property of film music rather than of this particular collection of 43 films. The
experiment in this section addresses that directly: a classifier is trained on the Eerola
corpus alone and applied, without any further fitting, to the 110 films of the Blockbuster
corpus of \citet{ma2021computational}. It also provides the first opportunity to compare
against an externally published result on the same data.

\subsection{Constructing a Shared Genre Space}
\label{subsec:shared_genres}

A model can only be transferred between two corpora if it predicts the same classes in
both. The two corpora do not annotate the same genre inventory: the Eerola target used up
to this point comprises eight genres including Crime, Adventure, Biography and
Documentary, none of which exist in the Blockbuster annotation, while Blockbuster includes
Sci-Fi and Romance, which are not among the eight. Both corpora are therefore relabelled
into a single shared space of six genres --- Action, Drama, Comedy, Sci-Fi, Romance and
Horror.

These six are not chosen here but adopted from \citet{ma2021computational}, and the
mapping they used to obtain them was recovered from their published film--genre master
list rather than reconstructed by assumption. It is the identity mapping: a film carries
reduced genre $g$ if and only if $g$ occurs in its raw IMDb genre list, which holds for
all 110 of their films without exception. Since this coincides with the ``any-present''
inclusion rule already applied to the Eerola annotations
(Section~\ref{subsec:genre_labels}), both corpora enter the shared space under the same
rule and without any corpus-specific heuristic. Applying it to Eerola retains 319 of the
360 clips, drawn from 41 films; the resulting label distributions are given in
Table~\ref{tab:shared_space}.

\begin{table}[t]
  \centering
  \caption{The shared six-genre label space. Both corpora are relabelled under the same
    ``any-present'' rule, so the same classifier can be trained on one and applied to the
    other.}
  \label{tab:shared_space}
  \begin{tabular}{lrr}
    \hline
    Genre   & Eerola (clips) & Blockbuster (films) \\
    \hline
    Action  & 103 & 55 \\
    Drama   & 239 & 44 \\
    Comedy  &  40 & 37 \\
    Sci-Fi  &  19 & 36 \\
    Romance &  51 & 13 \\
    Horror  &  28 & 11 \\
    \hline
    Total   & 319 clips (41 films) & 110 films \\
    \hline
  \end{tabular}
\end{table}

VGGish is the one representation both corpora provide identically --- extracted here for
the Eerola clips, published by \citeauthor{ma2021computational} for the Blockbuster films
--- and the two feature spaces are numerically compatible, both being the standard
post-processed 128-dimensional VGGish embedding with per-dimension means correlating at
$r=0.97$. A model fitted in one space can therefore be evaluated in the other.

\subsection{Reproducing the Published Result}

A transfer score cannot be interpreted without knowing what in-domain performance looks
like on the target corpus, and in particular without establishing that any shortfall is
attributable to the transfer rather than to a weak classifier. The protocol of
\citet{ma2021computational} --- leave-one-film-out cross-validation over the 110 films,
with precision, recall and $F_1$ computed per genre and macro-averaged --- was therefore
reproduced using the simple pipeline of this thesis: mean-pooled VGGish features and a
binary-relevance logistic regression. The result, macro-$F_1 = 0.620$, falls within the
range of $0.61$--$0.65$ that \citeauthor{ma2021computational} report for their VGGish
models, and both of their baselines reproduce to within $0.03$
(Table~\ref{tab:blockbuster_reference}). The pipeline is thus validated against an
external published result before being asked to transfer.

\begin{table}[t]
  \centering
  \caption{In-domain reference performance on the Blockbuster corpus
    (leave-one-film-out, macro-averaged). Published values are quoted from
    \citet{ma2021computational}, Table~4.}
  \label{tab:blockbuster_reference}
  \begin{tabular}{lrrr}
    \hline
    Model & Prec. & Rec. & $F_1$ \\
    \hline
    \citeauthor{ma2021computational}: VGGish, single-attention (their best) & .60 & .73 & .65 \\
    \citeauthor{ma2021computational}: VGGish, average pooling               & .55 & .78 & .62 \\
    \citeauthor{ma2021computational}: VGGish, $k$NN Simple-MI               & .64 & .59 & .61 \\
    \citeauthor{ma2021computational}: MIR, SVM Simple-MI                    & .60 & .52 & .56 \\
    \textbf{This thesis: VGGish + logistic regression}          & .569 & .705 & \textbf{.620} \\
    \hline
    Random-guess baseline (published / reproduced)              & .32 / .294 & .32 / .293 & .32 / .292 \\
    Plurality-label baseline (published / reproduced)           & .14 / .138 & .34 / .333 & .19 / .193 \\
    \hline
  \end{tabular}
\end{table}

The two baselines deserve a remark, because they are not interchangeable and the
literature uses both. The \emph{plurality} baseline always predicts the single most
frequent label set and reaches macro-$F_1 = 0.19$; the \emph{random-guess} baseline
predicts each genre independently at its base rate and reaches $0.32$, since the expected
macro-$F_1$ of such a predictor is approximately the mean base rate. Reporting a result
against the weaker of the two would overstate its margin by more than a tenth of an $F_1$
point, so both are computed throughout.

\subsection{Zero-Shot Transfer}

Every model in Table~\ref{tab:zero_shot} is fitted on the Eerola corpus alone and used to
predict all 110 Blockbuster films; no Blockbuster label participates in training. Where
the genre classifier consumes predicted emotions, those predictions are generated
out-of-fold on the training corpus, so that the classifier is fitted on the same noisy
estimates it will encounter at test time rather than on in-sample predictions that would
be unrealistically accurate. Because a single train--test split provides no fold structure
over which to form a confidence interval, the margins are quantified by a paired bootstrap
over the 110 target films, with both arms re-scored on each of 2000 resamples.

\begin{table}[t]
  \centering
  \caption{Zero-shot transfer: trained on Eerola, evaluated on all 110 Blockbuster films.
    The in-domain reference from Table~\ref{tab:blockbuster_reference} is repeated for
    orientation.}
  \label{tab:zero_shot}
  \begin{tabular}{lrrr}
    \hline
    Approach (trained on Eerola only) & Prec. & Rec. & $F_1$ \\
    \hline
    VGGish $\rightarrow$ predicted emotion (11), \textbf{per cue} & .488 & .651 & \textbf{.511} \\
    \quad the same, applied at film level                           & .450 & .637 & .482 \\
    Emotion (11) $\rightarrow$ genre, classifier fitted on ratings  & .445 & .474 & .434 \\
    PCA-8(VGGish) $\rightarrow$ genre (bottleneck control)          & .534 & .467 & .421 \\
    VGGish-128 $\rightarrow$ genre (direct)                         & .576 & .442 & .407 \\
    Random-guess baseline                                           & .294 & .293 & .292 \\
    \hline
    \emph{In-domain reference (trained on Blockbuster)}             & \emph{.569} & \emph{.705} & \emph{.620} \\
    \hline
  \end{tabular}
\end{table}

\begin{table}[t]
  \centering
  \caption{Paired bootstrap over the 110 target films (2000 resamples), macro-$F_1$
    differences under zero-shot transfer.}
  \label{tab:zero_shot_bootstrap}
  \begin{tabular}{lrrr}
    \hline
    Comparison & Diff. & 95\,\% CI & $p$ \\
    \hline
    Predicted emotion vs.\ VGGish direct & $+0.106$ & $[+0.015, +0.186]$ & $\mathbf{0.018}$ \\
    Predicted emotion vs.\ PCA-8 control & $+0.092$ & $[+0.009, +0.171]$ & $\mathbf{0.032}$ \\
    Film-level emotion vs.\ VGGish direct & $+0.077$ & $[-0.007, +0.152]$ & $0.066$ \\
    PCA-8 control vs.\ VGGish direct     & $+0.014$ & $[-0.032, +0.062]$ & $0.572$ \\
    \hline
  \end{tabular}
\end{table}

Four observations follow.

First, genre information transfers across the corpora at all. The best transferred model
reaches macro-$F_1 = 0.482$ against a random-guess floor of $0.292$, recovering
approximately $78\,\%$ of the $0.620$ obtained with in-domain training --- and it does so
despite a substantial change in the unit of analysis, since each training row is a
single ten-second excerpt while each test row aggregates the cues of an entire film. What
the models have learned is therefore a property of film music generally, not of the
43-film Eerola collection.

Second, and centrally for RQ2b, \emph{the emotion bottleneck transfers better than the
representation it is derived from}. The predicted-emotion pipeline exceeds the direct
VGGish classifier by $+0.106$ macro-$F_1$ ($p = 0.018$), a margin roughly twice the
$+0.04$--$0.05$ observed in-domain (Section~\ref{sec:statistical_power}). All figures use
the conservative of the two feature-scaling regimes, namely the one in which the baseline
performs better; under the alternative the same margin is $+0.116$ ($p<0.001$). This is the setting
in which an interpretable intermediate should be expected to help most: an emotion value
denotes the same perceptual quantity in both corpora, whereas an arbitrary embedding
dimension does not, and a classifier trained on the latter partly learns the idiosyncratic
scaling of its training corpus.

Third, this interpretation is supported by an ablation of the feature scaling.
Re-standardising each corpus using its own statistics --- an unsupervised adaptation that
uses no labels --- changes the direct embedding ($0.407 \rightarrow 0.394$) and the PCA
control ($0.421 \rightarrow 0.413$), but leaves the emotion pipeline unchanged at
$0.482$. Because emotion values are expressed in rating units rather than in embedding
units, the domain shift is absorbed by the representation itself.

Fourth, the advantage is not merely an effect of dimensionality reduction. Compressing the
same embedding to eight principal components yields only $+0.014$ over the raw features
($p = 0.572$), whereas compressing it to eight \emph{emotions} yields $+0.106$. The direct
comparison between the two bottlenecks favours emotion by $+0.092$
($p = 0.032$). Every in-domain version of this comparison was borderline at best; under
transfer it is unambiguous. It is therefore the \emph{emotion} bottleneck that helps, not
merely \emph{a} bottleneck --- which is the specific claim RQ2b makes.

Between the two granularities lies a third observation. Predicting emotion for each cue
and averaging the resulting emotion vectors outperforms averaging the embedding first and
predicting once ($0.511$ vs $0.482$; in-domain the same change is worth $+0.059$,
$p = 0.020$). A cue is a single piece of film music of median 19\,s, i.e.\ the same kind of
object as an Eerola excerpt, whereas a film-level mean over some 39 cues is not. Applying
the emotion model at the granularity it was trained on is thus not a detail but a
measurable part of the method.

\subsection{Do the Emotion--Genre Signatures Replicate?}
\label{subsec:signature_replication}

The per-genre emotional signatures of Section~\ref{subsec:emotion_genre} were derived on
the Eerola corpus from human ratings. Whether they describe film music or merely that
corpus is a question the second dataset can answer. The Eerola-trained regressor predicts
emotions for each of the 4664 Blockbuster cues, and the same effect sizes are recomputed
there; since no Blockbuster annotation exists or is used, any agreement is transfer rather
than fitting.

\begin{table}[t]
  \centering
  \caption{Replication of the per-genre emotional signatures on the Blockbuster corpus.
    $r(d)$ is the correlation between the two corpora's eight Cohen's $d$ values for that
    genre; sign agreement counts the emotions whose direction matches.}
  \label{tab:signature_replication}
  \begin{tabular}{lrrll}
    \hline
    Genre & $r(d)$ & Sign agr. & Eerola top & Blockbuster top \
    \hline
    Romance & $+0.99$ & 100\,\% & $+$valence & $-$anger \
    Action  & $+0.97$ &  88\,\% & $-$happy   & $+$anger \
    Sci-Fi  & $+0.97$ & 100\,\% & $+$anger   & $+$anger \
    Horror  & $+0.95$ &  88\,\% & $+$fear    & $+$fear \
    Drama   & $+0.80$ & 100\,\% & $+$valence & $-$energy \
    Comedy  & $+0.59$ &  75\,\% & $-$fear    & $+$happy \
    \hline
    \textbf{Pooled (48 cells)} & $\mathbf{+0.84}$ & \textbf{92\,\%} & & \
    \hline
  \end{tabular}
\end{table}

The signatures replicate. Across all 48 genre-by-emotion combinations the two corpora's
effect sizes correlate at $r = 0.84$ with $92\,\%$ sign agreement, and horror is defined by
fear in both. Where the single strongest emotion differs, it is a near-tie within a
signature of the same overall shape: action is low-valence, high-anger and low-happiness in
both corpora, and only the ranking of anger against happiness differs.

One caveat concerns magnitude rather than pattern. Measured per film, the Blockbuster
effect sizes are roughly twice the Eerola ones, because averaging some 39 cues per film
shrinks within-group variance and thereby inflates $d$. The cue-level estimates, computed
on the unit comparable to an Eerola excerpt, are of the same order as the originals
(horror--fear $+0.55$ against $+0.78$). What transfers is the structure of the
emotion--genre relationship, not the absolute effect size.

This is the strongest available evidence that the interpretability claim of this thesis
generalises beyond its own corpus: the affective fingerprints are a property of film
scoring practice rather than of one listener panel.

Finally, the magnitude of the drop from in-domain to transferred performance
($0.620 \rightarrow 0.482$) is itself consistent with the music literature.
\citet{eerola2011genrespecific} found that models of musical emotion generalise
considerably worse across musical repertoires than within them, explaining $16\,\%$ of
valence variance across genres against $43\,\%$ within. That the same asymmetry appears
one level up, in genre prediction, suggests it is a structural property of the
emotion--genre relationship rather than an artefact of either corpus.
